% Part: incompleteness
% Chapter: incompleteness-provability
% Section: rosser-thm

\documentclass[../../../include/open-logic-section]{subfiles}

\begin{document}

\olfileid{inc}{inp}{ros}

\olsection{રોસરનું પ્રમેય}

``$\omega$-સુસંગત''ને
બદલે માત્ર ``સુસંગત''
રાખીને વધુ મજબૂત
પરિણામ મેળવવા ગોડેલની
સાબિતી બદલી શકીએ?
રોસરે શોધેલી યુક્તિથી
જવાબ ``હા'' છે.
રોસરની યુક્તિ
$\OProv[\Th{T}](y)$ને
બદલે ``બદલાયેલું''
!!{derivability} વિધેયધર્મ
$\ORProv_T(y)$ વાપરે છે.

\begin{thm}
\ollabel{thm:rosser}
$\Th{T}$ એ~$\Th{Q}$નું
કોઈપણ સુસંગત,
!!{axiomatizable} વિસ્તરણ
હોય. તો $\Th{T}$
પૂર્ણ નથી.
\end{thm}

\begin{proof}
યાદ રાખો કે
$\OProv[\Th{T}](y)$ની
વ્યાખ્યા
$\lexists[x][\OPrf[\Th{T}](x,
  y)]$ છે; અહીં
$\OPrf[\Th{T}](x, y)$
એ એવો નિર્ણેય સંબંધ
નિરૂપે છે કે $x$ એ
ગોડેલ સંખ્યા~$y$
ધરાવતા !!{sentence}ના
!!a{derivation}ની ગોડેલ
સંખ્યા હોય ત્યારે અને
માત્ર ત્યારે જ સંબંધ
સાચો છે. $x$ અને~$y$
વચ્ચેનો એવો સંબંધ પણ
નિર્ણેય છે જેમાં $x$ એ
ગોડેલ સંખ્યા~$y$ ધરાવતા
વાક્યના \emph{ખંડન}ની
ગોડેલ સંખ્યા હોય ત્યારે
તે સાચો થાય. $\fn{not}(x)$
નીચે મુજબનું આદિમ
પુનરાવર્તી વિધેય હોય:
$x$ જો સૂત્ર~$!A$નો
સંકેત હોય, તો
$\fn{not}(x)$ એ
$\lnot
!A$નો સંકેત છે.
ત્યારે $\Refut[\Th{T}](x, y)$
ત્યારે અને માત્ર ત્યારે જ
સાચો છે જ્યારે
$\Prf[\Th{T}](x, \fn{not}(y))$
સાચો હોય. તેને
$\ORefut[\Th{T}](x, y)$
નિરૂપે એમ માનો. પછી
જો $\Th{T} \Proves \lnot !A$
અને $\delta$ એ અનુરૂપ
!!{derivation} હોય, તો
$\Th{Q} \Proves
\ORefut[\Th{T}](\gn{\delta}, \gn{!A})$.
$\ORProv[\Th{T}](y)$ની
વ્યાખ્યા આ રાખીએ:
\[
\lexists[x][(\OPrf[\Th{T}](x,y) \land \lforall[z][(z < x \lif \lnot
  \ORefut[\Th{T}](z,y))])].
\]
આશરે, $\ORProv[\Th{T}](y)$
કહે છે કે ``$\Th{T}$માં
$y$ની સાબિતી છે, અને
$y$નું તેની સાબિતીની
સંકેતસંખ્યા કરતાં નાની
સંકેતસંખ્યા ધરાવતું
ખંડન નથી.''
\sourcecorrection{OLINC-043}{સ્થિર અંગ્રેજી
મૂળમાં ``ટૂંકું ખંડન''
કહ્યું છે, પરંતુ વ્યાખ્યાનું
ક્રમ-માપ સાબિતી અને
ખંડનની સંકેતસંખ્યાઓની
સરખામણી છે. સૂત્રમાં
લખાયેલી નાની સંકેતસંખ્યા
અનુસાર અર્થ સ્પષ્ટ કર્યો છે.}
$\Th{T}$ સુસંગત
છે એમ માનીએ તો
$\ORProv[\Th{T}](y)$ અને
$\OProv[\Th{T}](y)$
એક જ સંખ્યાઓ માટે
સાચાં છે; પરંતુ
~$\Th{T}$માં
\emph{સાબિતિક્ષમતા}ની
દૃષ્ટિએ (સત્ય અને
સાબિતિક્ષમતા વચ્ચે તફાવત
છે તે હવે જાણીએ છીએ!)
બંનેના ગુણધર્મો જુદા છે.
જો $\Th{T}$
\emph{અ}સુસંગત હોય,
તો બંને એક જ સંખ્યાઓ
માટે સાચાં \emph{નથી}!{}
($\ORProv[\Th{T}](y)$ને
ઘણી વાર ``$y$ રોસર-સાબિતિક્ષમ
છે'' એમ વાંચે છે.
હમણાં ચર્ચ્યું તેમ,
રોસર-સાબિતિક્ષમતા કોઈ
ખાસ પ્રકારની સાબિતિક્ષમતા
નથી: અસુસંગત સિદ્ધાંતોમાં
કેટલાંક !!{sentence}
સાબિતિક્ષમ હોય છતાં
રોસર-સાબિતિક્ષમ નથી.
એટલે આ નામ ગૂંચવી શકે.
ગૂંચવણ ટાળવા તેને
``$y$ શ્મૂવેબલ છે'' પણ
વાંચી શકો.)

નિશ્ચિત-બિંદુ લેમ્માથી એવું
સૂત્ર $!R_\Th{T}$ છે કે
\begin{equation}
  \Th{Q} \Proves !R_\Th{T} \liff \lnot \ORProv[\Th{T}](\gn{!R_\Th{T}}).
  \ollabel{RT}
\end{equation}
\olref[1in]{thm:first-incompleteness}ની
સાબિતીથી વિપરીત, અહીં
દાવો એવો છે કે
$\Th{T}$ સુસંગત હોય તો
$\Th{T}$ એ $!R_\Th{T}$ને
!!{derive} કરતું નથી,
અને $\Th{T}$ એ
$\lnot !R_\Th{T}$ને પણ
!!{derive} કરતું નથી.
(બીજા શબ્દોમાં,
$\omega$-સુસંગતતાની
ધારણા જરૂરી નથી.)

પ્રથમ બતાવીએ કે
$\Th{T} \Proves/ !R_{\Th{T}}$.
માનો કે તે સાબિત કરે;
તો~$T$માંથી
~$!R_{\Th{T}}$નું
!!a{derivation} છે;
તેની ગોડેલ સંખ્યા~$n$
માનો. તેથી
$\Th{Q} \Proves \OPrf[\Th{T}](\num{n}, \gn{!R_{\Th{T}}})$,
કારણ કે $\OPrf[\Th{T}]$
એ~$\Th{Q}$માં
$\Prf[\Th{T}]$નું નિરૂપણ
કરે છે. ઉપરાંત દરેક
$k < n$ માટે $k$ એ
$\lnot !R_{\Th{T}}$ના
!!a{derivation}ની ગોડેલ
સંખ્યા નથી, કારણ કે
$\Th{T}$ સુસંગત છે.
તેથી દરેક $k < n$ માટે
$\Th{Q} \Proves \lnot
\ORefut[\Th{T}](\num{k}, \gn{!R_{\Th{T}}})$.
\olref[req][min]{lem:less-nsucc}
મુજબ
$\Th{Q} \Proves \lforall[z][(z < \num{n} \lif \lnot \ORefut[\Th{T}](z,
  \gn{!R_{\Th{T}}}))]$.
આથી
\[
\Th{Q} \Proves \lexists[x][(\OPrf[\Th{T}](x,\gn{!R_{\Th{T}}}) \land \lforall[z][(z
    < x \lif \lnot \ORefut[\Th{T}](z,\gn{!R_{\Th{T}}}))])],
\]
પણ આ તો
$\ORProv[\Th{T}](\gn{!R_{\Th{T}}})$
જ છે. \olref{RT} મુજબ
$\Th{Q}
\Proves \lnot !R_{\Th{T}}$.
$\Th{T}$ એ $\Th{Q}$નું
વિસ્તરણ હોવાથી
$\Th{T}
\Proves \lnot !R_{\Th{T}}$
પણ. પરંતુ
$\Th{T} \Proves !R_{\Th{T}}$
એવી ધારણા કરી હતી;
તેથી $\Th{T}$ અસુસંગત
બને, જે પ્રમેયની
ધારણાથી વિરુદ્ધ છે.

હવે બતાવીએ કે
$\Th{T} \Proves/ \lnot !R_{\Th{T}}$.
ફરી માનો કે તે
સાબિત કરે અને
~$\lnot !R_{\Th{T}}$ના
!!a{derivation}ની ગોડેલ
સંખ્યા~$n$ હોય. ત્યારે
$\Refut[\Th{T}](n, \Gn{!R_{\Th{T}}})$
સાચો છે; અને
$\ORefut[\Th{T}]$ એ
$\Th{Q}$માં
$\Refut[\Th{T}]$નું
નિરૂપણ કરે છે, તેથી
$\Th{Q} \Proves
\ORefut[\Th{T}](\num{n}, \gn{!R_{\Th{T}}})$.
અહીં પણ $\Th{T}$
અસુસંગત બનશે એવું
બતાવીશું, કારણ કે તે
~$!R_{\Th{T}}$ને પણ
!!{derive} કરશે. કારણ કે
\begin{align*}
\Th{Q} & \Proves !R_{\Th{T}} \liff \lnot \ORProv[\Th{T}](\gn{!R_{\Th{T}}}),
\intertext{અને $\Th{T}$ એ~$\Th{Q}$નું વિસ્તરણ હોવાથી, આ બતાવવું પૂરતું છે કે}
\Th{Q} & \Proves \lnot \ORProv[\Th{T}](\gn{!R_{\Th{T}}}).
\end{align*}
!!{sentence} $\lnot
\ORProv[\Th{T}](\gn{!R_{\Th{T}}})$,
એટલે
\begin{align*}
  \lnot & \lexists[x][(\OPrf[\Th{T}](x,\gn{!R_{\Th{T}}}) \land
    \lforall[z][(z < x \lif \lnot \ORefut[\Th{T}](z,\gn{!R_{\Th{T}}}))])],
  \intertext{તર્કની દૃષ્ટિએ આને સમકક્ષ છે:}
  & \lforall[x][(\OPrf[\Th{T}](x,\gn{!R_{\Th{T}}}) \lif
    \lexists[z][(z < x \land \ORefut[\Th{T}](z,\gn{!R_{\Th{T}}}))])].
\end{align*}
~$\Th{Q}$ જે હકીકતો
!!{derive} કરે છે તે
વાપરીને અનૌપચારિક
તાર્કિક દલીલ કરીએ.
$x$ ગમે તે હોય અને
$\OPrf[\Th{T}](x,
\gn{!R_{\Th{T}}})$ માનો.
આપણે પહેલેથી જાણીએ
છીએ કે
$\Th{T} \Proves/ !R_{\Th{T}}$;
તેથી દરેક $k$ માટે
$\Th{Q} \Proves \lnot \OPrf[\Th{T}](\num{k}, \gn{!R_{\Th{T}}})$.
આથી દરેક $k$ માટે
$\eq/[x][\num{k}]$ મળે
છે. ખાસ કરીને
(a) $\eq/[x][\num{n}]$.
ઉપરાંત
$\lnot(\eq[x][\num{0}]
\lor \eq[x][\num{1}] \lor \dots \lor \eq[x][\num{n-1}])$
મળે છે; તેથી
\olref[req][min]{lem:less-nsucc}
મુજબ (b)
$\lnot(x < \num{n})$.
\olref[req][min]{lem:trichotomy}
મુજબ $\num{n} < x$.
અને
$\Th{Q} \Proves
\ORefut[\Th{T}](\num{n}, \gn{!R_{\Th{T}}})$
હોવાથી
$\num{n} < x \land
\ORefut[\Th{T}](\num{n}, \gn{!R_{\Th{T}}})$,
અને તેથી
$\lexists[z][(z < x
\land \ORefut[\Th{T}](z,\gn{!R_{\Th{T}}}))]$
મળે છે. $x$ ગમે તે
હોવાથી માંગેલું વિધાન
મળે છે:
\[
\lforall[x][(\OPrf[\Th{T}](x,\gn{!R_{\Th{T}}}) \lif
  \lexists[z][(z < x \land \ORefut[\Th{T}](z,\gn{!R_{\Th{T}}}))])].
\]
\end{proof}

\begin{prob}
પ્રાકૃતિક સંખ્યાઓના બે
ગણો $A$ અને $B$ને
\emph{સંગણનીય રીતે
અવિભાજ્ય} ત્યારે કહીએ
જ્યારે એવો કોઈ
!!{decidable} ગણ~$X$
ન હોય કે
$A \subseteq X$ અને
$B \subseteq \Complement{X}$
(અહીં $\Complement{X}$
એ~$X$નો પૂરક ગણ
$\Nat \setminus X$ છે).
$\Th{T}$ એ~$\Th{Q}$નું
સુસંગત,
!!{axiomatizable} વિસ્તરણ
હોય. માનો કે $A$ એ
~$\Th{T}$માં સાબિત
થતાં !!{sentence}ની
ગોડેલ સંખ્યાઓનો ગણ છે,
અને $B$ એ~$\Th{T}$માં
નકારી શકાય એવાં
વાક્યોની ગોડેલ
સંખ્યાઓનો ગણ છે.
$A$ અને~$B$ સંગણનીય
રીતે અવિભાજ્ય છે એમ
સાબિત કરો.
\end{prob}

\end{document}
